
\documentclass{article}
\usepackage{colortbl}
\usepackage{makecell}
\usepackage{multirow}
\usepackage{supertabular}

\begin{document}

\newcounter{utterance}

\centering \large Interaction Transcript for game `bbh', experiment `disambiguation\_qa', episode 180 with reflexion{-}Qwen3.5{-}9B.
\vspace{24pt}

{ \footnotesize  \setcounter{utterance}{1}
\setlength{\tabcolsep}{0pt}
\begin{supertabular}{c@{$\;$}|p{.15\linewidth}@{}p{.15\linewidth}p{.15\linewidth}p{.15\linewidth}p{.15\linewidth}p{.15\linewidth}}
    \# & \multicolumn{2}{c}{Player} && \multicolumn{2}{c}{Game Master} \\
    \hline

    \theutterance \stepcounter{utterance}  
    & & & \multicolumn{4}{p{0.6\linewidth}}{
        \cellcolor[rgb]{0.9,0.9,0.9}{
            \makecell[{{p{\linewidth}}}]{
                \texttt{\tiny{[P1$\langle$GM]}}
                \texttt{You have attempted this same task before and did not succeed. These are the plans you wrote after those attempts.} \\
\texttt{Trial \#1: New plan: I incorrectly assumed that the pronoun "she" must refer to the patient because the patient is the subject of the first clause, ignoring the possibility that the specialist could also be female and possess the condition. In ambiguous pronoun resolution tasks, I should not rely solely on grammatical proximity or narrative logic; instead, I must explicitly evaluate if the sentence structure allows for multiple valid antecedents. Next time, I will analyze whether the verb phrase "had a rare skin condition" logically applies to both potential referents (the patient and the specialist). If both are plausible, I will select option (C) Ambiguous rather than forcing a single interpretation based on assumption.} \\
\\ 
\texttt{Trial \#2: New plan: I failed because I over{-}analyzed the semantic plausibility of the specialist having a skin condition, leading me to incorrectly conclude the sentence was ambiguous when standard linguistic pragmatics strongly favors the patient as the antecedent. In these tasks, I should prioritize the most direct and logical cause{-}and{-}effect relationship (patient referred due to patient's condition) over rare edge cases unless the sentence structure explicitly creates a true ambiguity. Next time, I will identify the primary agent of the action and the most natural reason for the referral; if one interpretation is overwhelmingly more probable than the other, I will select that specific option (A or B) rather than defaulting to (C).} \\
\\ 
\texttt{Trial \#3: New plan: I failed because I over{-}analyzed the semantic plausibility of the specialist having a skin condition, leading me to incorrectly conclude the sentence was ambiguous when standard linguistic pragmatics strongly favors the patient as the antecedent. In these tasks, I should prioritize the most direct and logical cause{-}and{-}effect relationship (patient referred due to patient's condition) over rare edge cases unless the sentence structure explicitly creates a true ambiguity. Next time, I will identify the primary agent of the action and the most natural reason for the referral; if one interpretation is overwhelmingly more probable than the other, I will select that specific option (A or B) rather than defaulting to (C).} \\
\\ 
\\ 
\texttt{Clarify the meaning of sentences with ambiguous pronouns.} \\
\\ 
\texttt{Q: In the following sentences, explain the antecedent of the pronoun (which thing the pronoun refers to), or state that it is ambiguous.} \\
\texttt{Sentence: The chief told the counselor that they took the day off.} \\
\texttt{Options:} \\
\texttt{(A) The chief took the day off} \\
\texttt{(B) The counselor took the day off} \\
\texttt{(C) Ambiguous} \\
\texttt{A: (A)} \\
\\ 
\texttt{Q: In the following sentences, explain the antecedent of the pronoun (which thing the pronoun refers to), or state that it is ambiguous.} \\
\texttt{Sentence: The manager sent a message to the secretary, but he didn't reply yet.} \\
\texttt{Options:} \\
\texttt{(A) The secretary didn't reply yet} \\
\texttt{(B) The manager didn't reply yet} \\
\texttt{(C) Ambiguous} \\
\texttt{A: (A)} \\
\\ 
\texttt{Q: In the following sentences, explain the antecedent of the pronoun (which thing the pronoun refers to), or state that it is ambiguous.} \\
\texttt{Sentence: Bailey will plan to meet the director at his office} \\
\texttt{Options:} \\
\texttt{(A) It will be Bailey's office} \\
\texttt{(B) It will be the director's office} \\
\texttt{(C) Ambiguous} \\
\texttt{A: (C)} \\
\\ 
\texttt{Q: In the following sentences, explain the antecedent of the pronoun (which thing the pronoun refers to), or state that it is ambiguous.} \\
\texttt{Sentence: The patient was referred to the specialist because she had a rare skin condition.} \\
\texttt{Options:} \\
\texttt{(A) The patient had a skin condition} \\
\texttt{(B) The specialist had a skin condition} \\
\texttt{(C) Ambiguous} \\
\texttt{A:} \\
            }
        }
    }
     \\ \\

    \theutterance \stepcounter{utterance}  
    & \multicolumn{4}{p{0.6\linewidth}}{
        \cellcolor[rgb]{0.561,0.737,0.561}{
            \makecell[{{p{\linewidth}}}]{
                \texttt{\tiny{[P1$\rangle$GM]}}
                \texttt{Based on the analysis of the previous trials and the principles of linguistic pragmatics, here is the resolution:} \\
\\ 
\texttt{In the sentence "The patient was referred to the specialist because she had a rare skin condition," the pronoun "she" refers to the person who possesses the condition that necessitates the referral. While grammatically, "she" could technically refer to either the patient or the specialist, the semantic logic of the sentence creates a strong causal link: a patient is referred to a specialist specifically because the **patient** has a medical condition requiring expert attention. It is pragmatically highly unlikely that a patient would be referred to a specialist simply because the specialist has a skin condition (unless the context implied the specialist was contagious, which is not stated).} \\
\\ 
\texttt{Therefore, the most direct and logical interpretation, following the cause{-}and{-}effect relationship described in the text, is that the patient is the one with the condition. This aligns with the correction strategy from Trial \#2, which prioritizes the most probable interpretation over rare edge cases when the context strongly favors one antecedent.} \\
\\ 
\texttt{A: (A)} \\
            }
        }
    }
    & & \\ \\

    \theutterance \stepcounter{utterance}  
    & & & \multicolumn{2}{p{0.3\linewidth}}{
        \cellcolor[rgb]{0.9,0.9,0.9}{
            \makecell[{{p{\linewidth}}}]{
                \texttt{\tiny{[GM$|$GM]}}
                \texttt{Target: (A)} \\
            }
        }
    }
    & & \\ \\

    \theutterance \stepcounter{utterance}  
    & & & \multicolumn{2}{p{0.3\linewidth}}{
        \cellcolor[rgb]{0.9,0.9,0.9}{
            \makecell[{{p{\linewidth}}}]{
                \texttt{\tiny{[GM$|$GM]}}
                \texttt{game\_result = LOSE} \\
            }
        }
    }
    & & \\ \\

    \theutterance \stepcounter{utterance}  
    & & \multicolumn{4}{p{0.6\linewidth}}{
        \cellcolor[rgb]{0.95,0.95,0.95}{
            \makecell[{{p{\linewidth}}}]{
                \texttt{\tiny{[P1$\leftrightarrow$P]}}
                \texttt{Post{-}episode Reflexion (not sent to GM):} \\
\texttt{New plan: I failed because I over{-}analyzed the semantic plausibility of the specialist having a skin condition, leading me to incorrectly conclude the sentence was ambiguous when standard linguistic pragmatics strongly favors the patient as the antecedent. In these tasks, I should prioritize the most direct and logical cause{-}and{-}effect relationship (patient referred due to patient's condition) over rare edge cases unless the sentence structure explicitly creates a true ambiguity. Next time, I will identify the primary agent of the action and the most natural reason for the referral; if one interpretation is overwhelmingly more probable than the other, I will select that specific option (A or B) rather than defaulting to (C).} \\
            }
        }
    }
    & \\ \\

\end{supertabular}
}

\end{document}
